%==============================================================
%  Liyan Tan — Curriculum Vitae
%  Compile with:  pdflatex cv_liyan_tan.tex
%==============================================================
\documentclass[10pt,letterpaper]{article}

\usepackage[top=1.15cm,bottom=1.15cm,left=1.5cm,right=1.5cm,footskip=0.6cm]{geometry}
\usepackage{charter}
\usepackage[T1]{fontenc}
\usepackage[utf8]{inputenc}
\usepackage[dvipsnames]{xcolor}
\usepackage{enumitem}
\usepackage{array}
\usepackage{tabularx}
\usepackage[letterspace=90]{microtype}
\usepackage{textcomp}
\usepackage{needspace}
\usepackage[hidelinks]{hyperref}

\definecolor{ink}{RGB}{26,26,26}          % 正文
\definecolor{navy}{RGB}{31,58,95}         % 姓名 + 顶部粗线
\definecolor{accent}{RGB}{10,102,114}     % 深青：小标题 + 分隔线
\definecolor{muted}{RGB}{95,102,108}      % 日期、联系方式

\pagestyle{empty}
\setlength{\parindent}{0pt}
\setlength{\tabcolsep}{0pt}
\color{ink}
\linespread{1.06}\selectfont

% ---------- 小标题：加宽字距的小型大写 + 下方细线 ----------
\newcommand{\cvsection}[1]{%
  \vspace{0.95em}%
  \noindent{\large\bfseries\color{accent}\textls[110]{\textsc{#1}}}\par
  \vspace{0.12em}%
  {\color{accent}\hrule height 0.7pt}%
  \vspace{0.55em}%
}

% ---------- 左侧内容 + 右侧日期 ----------
\newcommand{\entry}[2]{%
  \noindent\begin{tabularx}{\textwidth}{@{}X@{\hspace{1.4em}}r@{}}
    #1 & {\color{muted}#2}
  \end{tabularx}\par
}

\newlist{cvitems}{itemize}{1}
\setlist[cvitems]{label=\textbullet, leftmargin=1.1em, labelsep=0.4em,
                  topsep=0.28em, itemsep=0.22em, parsep=0pt, partopsep=0pt}

\newcommand{\pub}[2]{%
  \vspace{0.55em}
  \noindent\textbf{#1}\par
  \noindent#2\par
}

\newcommand{\me}{\textbf{Liyan Tan}}
\newcommand{\paperlink}[1]{\,{\small\color{accent}\href{#1}{[paper]}}}

\begin{document}

%================= HEADER =================
\noindent{\fontsize{22}{24}\selectfont\bfseries\color{navy} Liyan Tan}\hfill
{\small\color{muted}
  +1 (805) 259-0248 \ \textperiodcentered\ \href{mailto:liyan_tan@ucsb.edu}{liyan\_tan@ucsb.edu}
  \ \textperiodcentered\ \href{https://liyantan111.github.io}{liyantan111.github.io}
  \ \textperiodcentered\ \href{https://scholar.google.com/citations?user=xBvnXWsAAAAJ}{Google Scholar}
}

\vspace{0.42em}
{\color{navy}\hrule height 1.6pt}
\vspace{0.5em}

\noindent Zeroth-order and memory-efficient optimization for large language model fine-tuning and analog/RF circuit design.

%================= EDUCATION =================
\cvsection{Education}

{\small
\entry{\textbf{University of California, Santa Barbara}, \textit{Ph.D. in Computer Engineering} \textbar{} CA, USA}{Spring 2025 - Expected 2028}
\entry{\textbf{University of California, Santa Barbara}, \textit{M.S. in Computer Engineering} \textbar{} CA, USA}{Fall 2023 - Winter 2025}
\entry{\textbf{University of California, Santa Barbara}, \textit{Extension Program in ECE} \textbar{} CA, USA}{Fall 2022 - Spring 2023}
\entry{\textbf{Huazhong University of Science \& Technology}, \textit{B.S. in Electronic Information Engineering} \textbar{} China}{Fall 2019 - Summer 2023}
}

%================= INDUSTRIAL EXPERIENCE =================
\cvsection{Industrial Experience}

\entry{\textbf{Cadence Design Systems,} \textit{Software Architect Intern} \textbar{} Austin, TX}{Jun 2025 - Sep 2025}
\begin{cvitems}
  \item Architected an LLM copilot agent for Voltus, Cadence's power-integrity signoff solver, turning natural-language design intent into verified command invocations, GUI actions, and automated root-cause analysis over signoff reports.
  \item Enabled on-premise deployment under enterprise data constraints through retrieval over EDA documentation, parameter-efficient fine-tuning, and distillation of a large teacher into a compact local model.
\end{cvitems}

%================= PUBLICATIONS =================
\cvsection{Selected Publications}

\pub{GRZO: Group-Relative Zeroth-Order Optimization for Large Language Model Fine-Tuning\paperlink{https://arxiv.org/abs/2606.02857}}
{\me, Yequan Zhao, Yifan Yang, Ruijie Zhang, Xinling Yu, Zheng Zhang, \textit{Findings of the Association for Computational Linguistics: EMNLP}, 2026.}

\pub{ZOAF: Towards Efficient Zeroth-Order Optimization for Analog/RF Circuit Design\paperlink{https://arxiv.org/abs/2606.02869}}
{\me, Yequan Zhao, Jinming Lu, Ben F. Jamroz, Ari Feldman, Zheng Zhang, under review at \textit{IEEE TCAD}, arXiv:2606.02869, 2026.}

\pub{IAPO: Input Attribution-Aware Policy Optimization for Tool Use in Small Multimodal Agents\paperlink{https://arxiv.org/abs/2606.11652}}
{Yifan Yang, Zhen Zhang, Jiayi Tian, \me, Zheng Zhang, \textit{arXiv preprint arXiv:2606.11652}, 2026.}

\pub{FuRA: Full-Rank Parameter-Efficient Fine-Tuning with Spectral Preconditioning\paperlink{https://arxiv.org/abs/2605.22869}}
{Yequan Zhao, Ruijie Zhang, \me, Niall Moran, Tong Qin, Zheng Zhang, \textit{arXiv preprint arXiv:2605.22869}, 2026.}

\pub{MUON+: Towards More Effective Muon via One Additional Normalization Step for LLM Pre-training\paperlink{https://arxiv.org/abs/2602.21545}}
{Ruijie Zhang, Yequan Zhao, Ziyue Liu, Zhengyang Wang, Yupeng Su, \me, Zheng Zhang, \textit{arXiv preprint arXiv:2602.21545}, 2026.}

\pub{Overview and Idea of the Mars UAV Experimental Platform Development\paperlink{https://scholar.google.com/citations?view_op=view_citation&hl=en&user=xBvnXWsAAAAJ&citation_for_view=xBvnXWsAAAAJ:UeHWp8X0CEIC}}
{\me, Yijun Mao, \textit{Advances in Aeronautical Science and Engineering}, 15(2):1-10, 2024.}

%================= RESEARCH PROJECTS =================
\cvsection{Research Projects}

\needspace{5\baselineskip}
\entry{\textbf{Variance-Reduced Zeroth-Order Optimization for LLM Fine-Tuning} \textbar{} Advisor: Zheng Zhang}{Mar 2025 - Present}
\begin{cvitems}
  \item Proposed GRZO, a group-relative zeroth-order optimizer that assigns a pseudo-independent perturbation to every mini-batch example, raising the effective gradient-direction count from one to the batch size at no additional forward cost.
  \item Proved directional unbiasedness with variance shrinking proportionally to the batch size; at the same inference-level memory footprint and comparable wall-clock cost, improved Llama3-8B average accuracy by $+3.0$ over MeZO and lifted sparse, low-rank, and quantized ZO variants by $+4.9$ on average.
\end{cvitems}

\vspace{0.6em}
\needspace{5\baselineskip}
\entry{\textbf{Zeroth-Order Optimization for Analog/RF Circuit Design} \textbar{} Advisor: Zheng Zhang, with NIST}{Jan 2025 - Present}
\begin{cvitems}
  \item Proposed ZOAF, a surrogate-free framework recovering gradient-descent directions from a small number of black-box circuit simulations, with a hybrid schedule from random-direction exploration to coordinate-wise refinement, quasi-random multi-start, and a sliding-window monitor for early stopping and feasibility.
  \item Achieved the best median figure of merit on all three benchmark schematics, with up to an order-of-magnitude advantage in median peaking on a 22-parameter amplifier and $1.3$-$3.8\times$ fewer simulator calls to convergence.
  \item Extending the framework toward uncertainty-aware design: stochastic zeroth-order optimization to handle process variations in nanoscale circuits, and distributionally robust circuit optimization that accounts for distribution shift in those variations.
\end{cvitems}

\vspace{0.6em}
\needspace{5\baselineskip}
\entry{\textbf{Tensor-Train Compressed NN via Temporal Logic} \textbar{} Advisors: Zheng Zhang, Tim Sherwood}{Apr 2024 - Jan 2025}
\begin{cvitems}
  \item Implemented tensor-train compressed neural networks on temporal-logic hardware using delay-space values and approximated arithmetic, reducing parameter count by $115\times$ at 2.21\% (exact) and 4.72\% (approximated) accuracy loss.
  \item Cut latency by $14.28\times$ and energy by $32.37\times$ on a $4\times4$ PE array under race-logic hardware settings.
\end{cvitems}

\vspace{0.6em}
\needspace{5\baselineskip}
\entry{\textbf{Mars UAV Aerodynamic Performance Experimental Platform} \textbar{} Advisor: Yijun Mao, HUST}{Jan 2021 - Sep 2022}
\begin{cvitems}
  \item Surveyed two decades of Mars UAV experimental platforms and summarized their design trade-offs; the resulting review was published in \textit{Advances in Aeronautical Science and Engineering}.
  \item Built a hover-efficiency test platform with an axially distributed model, measuring rotational speed, torque, lift, and induced velocity of the rotor blades.
\end{cvitems}

%================= TEACHING =================
\cvsection{Teaching}

\noindent\textbf{Teaching Assistant, UC Santa Barbara:} PHYS 127BL Digital Electronics; PHYS 127AL Analog Electronics; PHYS 6BL Introductory Physics; ECE 15A Fundamentals of Logic Design; PSTAT 5A Understanding Data.

%================= HONORS =================
\cvsection{Honors}

\entry{\textbf{Honours Degree for Undergraduates}, Huazhong University of Science and Technology}{Jun 2023}

\end{document}
